\documentclass[12pt,a4paper]{report}

% 1. ตั้งค่าระยะขอบกระดาษตาม PDF: Left 1.5", Right 1", Top 1", Bottom 1"
\usepackage[left=1.5in, right=1in, top=1in, bottom=1in]{geometry}

% 2. ตั้งค่าระบบภาษาและฟอนต์ (ต้อง Compile ด้วย XeLaTeX)
\usepackage{fontspec}
\usepackage{polyglossia}
\setdefaultlanguage{english}
\setotherlanguage{thai}
\XeTeXlinebreaklocale "th"
\XeTeXlinebreakskip = 0pt plus 1pt

% 3. ตั้งค่าระยะห่างบรรทัด
\usepackage{setspace}

% 4. นำเข้าฟอนต์ Angsana New (ตามที่ระบุใน PDF)
\setmainfont[
    Path=fonts/,
    Extension=.ttf,
    BoldFont=AngsanaNew_Bold,
    ItalicFont=AngsanaNew_Italic,
    BoldItalicFont=AngsanaNew_BoldItalic,
]{AngsanaNew}

\begin{document}

% -------------------------
% หน้าปก (Cover Page)
% -------------------------
\begin{titlepage}
\begin{center}
    % ชื่อเรื่อง (20 pt Bold, ระยะห่างบรรทัด 1.5)
    \setstretch{1.5} 
    {\fontsize{20pt}{24pt}\selectfont \textbf{ASSESSMENT OF THE PERSPECTIVES OF PRACTITIONERS ON\\CARBON TAX MECHANISM TO ENCOURAGE THE USE OF ELECTRIC\\VEHICLES USED IN ROAD FREIGHT TRANSPORT ACTIVITIES IN\\THAILAND}}\\
    
    \vspace{2.5cm} % ปรับระยะห่างระหว่างชื่อเรื่องกับชื่อผู้เขียน
    
    % ชื่อผู้จัดทำ (อยู่กึ่งกลางหน้ากระดาษ)
    {\fontsize{16pt}{20pt}\selectfont \textbf{AUTHOR NAME IN ENGLISH}}\\
    
    \vfill
    \setstretch{1.0} % กลับมาระยะห่างบรรทัดปกติ 1.0
    
    % ส่วนท้ายของหน้าปก
    {\fontsize{16pt}{20pt}\selectfont \textbf{A THESIS REPORT SUBMITTED IN PARTIAL FULFILLMENT\\OF THE REQUIREMENTS FOR THE DEGREE OF\\MASTER OF SCIENCE IN INFORMATION TECHNOLOGY\\SCHOOL OF INFORMATION TECHNOLOGY\\KING MONGKUT'S INSTITUTE OF TECHNOLOGY LADKRABANG\\2021}}\\ % แก้ไขปีให้เป็นปีปัจจุบัน
    
    \vspace{1cm}
    
    % รหัสวิทยานิพนธ์
    {\fontsize{16pt}{20pt}\selectfont \textbf{KMITL-2025-IT-M-01-001}}
\end{center}
\end{titlepage}

% -------------------------
% บทคัดย่อ (Abstract)
% -------------------------
\chapter*{ABSTRACT}
\addcontentsline{toc}{chapter}{ABSTRACT}
Write your abstract here...

% -------------------------
% สารบัญ
% -------------------------
\tableofcontents
\clearpage

% -------------------------
% บทที่ 1 (Chapter 1)
% -------------------------
\chapter{INTRODUCTION}
This is the first chapter. The main content will be typed in Angsana New font as specified in the thesis format.

\vfill
\clearpage

% -------------------------
% ภาคผนวก (Appendix)
% -------------------------
\appendix
\chapter{APPENDIX A TITLE}
Appendix A content goes here.

% -------------------------
% ประวัติผู้จัดทำ (Author Biography) ตามหน้าสุดท้ายของ PDF
% -------------------------
\chapter*{AUTHOR BIOGRAPHY}
\addcontentsline{toc}{chapter}{AUTHOR BIOGRAPHY}

\noindent \textbf{Author:} \hspace{2cm} Name Surname \\
\textbf{Degree:} \hspace{1.9cm} Master of Science in Information Technology \\
\textbf{Date:} \hspace{2.2cm} DD/MM/YYYY \\
\textbf{Date of Birth:} \hspace{1cm} DD/MM/YYYY \\
\textbf{Place of Birth:} \hspace{0.9cm} Province, Country \\

\vspace{12pt}
\noindent \textbf{Undergraduate and Graduate Education:}
\begin{itemize}
    \setlength\itemsep{0em} % ระยะห่างระหว่าบรรทัดในลิสต์ตาม PDF
    \item Master of Science, King Mongkut's Institute of Technology Ladkrabang, Bangkok, 2017
    \item Bachelor degree in Business Management, University of Cambodia, Phnom Penh, 2014\\
    Major: Logistics and Supply Chain Management
\end{itemize}

\vspace{12pt}
\noindent \textbf{Presentations and Publications:}
\begin{itemize}
    \setlength\itemsep{0em}
    \item Your publication 1...
    \item Your publication 2...
\end{itemize}

\end{document}